\chapter{Working with MSA lanes}
\MARSChapterHeading{Working with MSA lanes}
\label{revision:msa}

A pixel conversion applies the same small calculation to many independent pixels.
That makes it a useful first problem for the MIPS SIMD Architecture, or MSA:
one instruction can perform the same operation on several elements of a vector
register. The difficult part is often arranging the elements so that those
parallel operations mean the calculation we intended.

This chapter converts a tightly packed array of red, green and blue bytes to
sixteen-bit pixels. You already know scalar loads, bit fields, array bounds and
function interfaces. We will keep those obligations while changing the work done
by the loop body. The complete public \texttt{rgb-formats.asm} in
\texttt{revision/labs/files} implements original project \texttt{019-rgb-formats}.
It includes packing, expansion and a two-row preview. Here we develop its packing
kernel; the expansion function provides a runnable comparison rather than another
construction sequence. The independent task changes the packed representation.

\section{Specify the pixel before choosing instructions}

RGB888 stores each pixel as three consecutive unsigned bytes: red $R$, green $G$
and blue $B$. The next pixel begins immediately afterward. There is no alpha byte
or padding. RGB565 stores one numerical halfword with five red bits, six green
bits and five blue bits. Packing discards low component bits:
\[
 r=\lfloor R/8\rfloor,\qquad
 g=\lfloor G/4\rfloor,\qquad b=\lfloor B/8\rfloor,
 \qquad P=2^{11}r+2^5g+b.
\]
The fields occupy disjoint bit positions, so bitwise OR combines them just as
addition does here. This is truncation, not rounding to the nearest value.

For $(R,G,B)=(123,45,67)$, the reduced fields are $(15,11,8)$. Their placed values
are $30720$, $352$ and $8$, giving $P=31080$, hexadecimal \texttt{7968}.
The required byte order writes \texttt{68 79}. A halfword's numerical value and
its two stored bytes are different descriptions of the same result.

\begin{figure}[htbp]
\centering
\begin{tikzpicture}[x=5.7mm,y=7mm,alt={An RGB565 halfword has bits fifteen through eleven for red, ten through five for green, and four through zero for blue. The example fields are red fifteen, green eleven and blue eight. Together they form hexadecimal 7968, stored as low byte 68 followed by high byte 79.}]
\fill[black!12] (0,0) rectangle (5,1);
\fill[black!22] (5,0) rectangle (11,1);
\fill[black!6] (11,0) rectangle (16,1);
\draw (0,0) rectangle (16,1) (5,0)--(5,1) (11,0)--(11,1);
\node at (2.5,.5) {Red: 15};\node at (8,.5) {Green: 11};\node at (13.5,.5) {Blue: 8};
\node[above] at (0,1) {15};\node[above] at (5,1) {10};\node[above] at (11,1) {4};\node[above] at (16,1) {0};
\node[below,align=center] at (8,-.25) {Numerical halfword: \texttt{7968}\\Stored bytes, increasing address: \texttt{68 79}};
\end{tikzpicture}
\caption{Field positions describe a numerical value. The byte order determines which part of that value is stored first. Shading distinguishes the fields without assigning display colors to their bits.}
\label{fig:revision-msa-fields}
\end{figure}


The supplied expansion repeats high bits into the newly available low positions:
\[
 R'=8r+\lfloor r/4\rfloor,\quad
 G'=4g+\lfloor g/16\rfloor,\quad
 B'=8b+\lfloor b/4\rfloor.
\]
Our example expands to $(123,44,66)$. Maximum fields expand to 255. Discarded
input bits cannot generally be recovered: distinct RGB888 pixels can share one
RGB565 encoding. Nor must every expanded component be below its original; red
128 packs to 16 and expands to 132. These are code-value rules, not a color-space
or gamma conversion. A plausible-looking picture cannot establish exact compliance.

The function \texttt{pack\_rgb565} receives full source and destination pointers
in \texttt{\$a0/\$a1} and count $n$ in \texttt{\$a2}, from zero through 32. The
selected input contains $3n$ bytes and output $2n$ bytes. They are disjoint and
may start at any byte alignment, including inside larger allocations. Return $n$
in \texttt{\$v0}, preserve all bits of the saved registers, frame pointer and
stack pointer, and return through the supplied continuation. Only selected output
bytes may change. Source bytes and padding remain the caller's property even when
mapped. This leaf function needs no stack frame because it uses no saved registers
or nested call.

\section{Read one register through several lane widths}

MSA provides 32 vector registers, \texttt{\$w0} through \texttt{\$w31}, each holding
128 bits. An instruction's element suffix partitions those bits into lanes:
\texttt{.b} uses sixteen eight-bit lanes, \texttt{.h} eight sixteen-bit lanes,
\texttt{.w} four thirty-two-bit lanes and \texttt{.d} two sixty-four-bit lanes.
A register does not permanently become a byte or halfword array. The next
instruction can interpret its unchanged bits at a different width.

In this little-endian MARS environment, a byte load to a vector puts the lowest
addressed byte in byte lane zero. Byte lanes zero and one form halfword lane zero,
with lane zero as its low byte. Thus bytes \texttt{7b 00} are halfword 123; bytes
\texttt{7b 2d} are halfword \texttt{2d7b}, not two independent component values
when a halfword instruction reads them. Widening a component therefore requires
arranging a zero high byte, not merely changing the instruction suffix.

Signedness also belongs to the operation and interpretation. A component byte
\texttt{ff} means unsigned 255 in this interface. A scalar signed byte load would
produce $-1$; \texttt{lbu} produces 255. The vector path gathers each component into
a zero-extended halfword and uses logical right shifts. The suffix alone does not
make data signed or unsigned. MSA registers overlap floating-point storage in this
target; keeping an unrelated floating-point value live in that storage across a
vector update would overwrite it. The exercise permits these registers to change.

Unlike a scalar byte load, \texttt{ld.b} loads an entire sixteen-byte vector.
Its byte suffix describes element organization, not a one-byte transfer. The
matching \texttt{st.b} stores sixteen bytes. Byte-element vector transfers permit
arbitrary byte alignment in the selected environment. Ordinary scalar halfword
loads and stores also accept unaligned addresses on this Release 6 target; byte
operations are not an alignment workaround. Every transfer still needs its whole
range inside the caller's selected slice.

\section{Cover eight pixels with two bounded windows}

Eight RGB888 pixels occupy 24 bytes, while their packed outputs occupy sixteen.
The worked group loads sixteen bytes at source offsets zero and eight:
\begin{minted}{mars-idea}
    ld.b $w0, 0($a0)
    ld.b $w1, 8($a0)
\end{minted}
The windows cover bytes 0 through 15 and 8 through 23. Their eight-byte overlap
is deliberate: all needed input is available in two registers, and both loads
remain within one complete group. Loading at offsets zero and sixteen would read
through byte 31. The extra mapped bytes would still be outside a 24-byte group
when it is the final group. A wide load cannot be justified by ignoring its
unwanted lanes afterward.

Before either load, the loop proves at least eight pixels remain. No vector
instruction runs on seven pixels, even if the input allocation happens to have
padding. The proof has two parts: the group gate supplies 24 input bytes and the
output contract supplies sixteen writable bytes. It is not enough to prove only
the first address valid. We use full-width pointer additions for subsequent groups;
the pixel count and component shifts need only their bounded scalar widths.

\section{Derive a shuffle from the loaded bytes}

\texttt{vshf.b} selects each output byte using a control byte from its destination
register. For \texttt{vshf.b \$w4,\$w1,\$w0}, selectors 0 through 15 choose byte
lanes in \texttt{\$w0}; selectors 16 through 31 choose lanes in \texttt{\$w1}.
The rightmost source supplies the lower selector range. Selector 64 produces zero.
We use only these explicitly defined controls in the worked tables.

The two source ranges are a selector namespace, not consecutive source addresses:
selector 16 refers to \texttt{\$w1} lane zero, which came from source byte eight.
Red components lie at input offsets $0,3,6,9,12,15,18,21$. The first six are
available in \texttt{\$w0}. The last two are lanes ten and thirteen of
\texttt{\$w1}, so their selectors are 26 and 29. Insert zero between successive
selected red bytes to obtain eight unsigned halfwords:
\begin{minted}{mars-idea}
gather_red: .byte 0,64,3,64,6,64,9,64,12,64,15,64,26,64,29,64
\end{minted}
Each pair is a low component byte followed by a zero high byte. The first default
pixel therefore gives halfword 123, not a signed extension and not a red/green
pair. Derive the green and blue tables by locating their actual source bytes;
adding one to every red selector would fail where the chosen window changes.
Green pixel five, at source offset sixteen, uses selector 24 rather than 16.

\begin{table}[htbp]
  \centering\small
  \begin{tabular}{lrrrrrrrr}
    \toprule
    Pixel              & 0 & 1 & 2 & 3  & 4  & 5  & 6  & 7  \\
    \midrule
    Red input offset   & 0 & 3 & 6 & 9  & 12 & 15 & 18 & 21 \\
    Red selector       & 0 & 3 & 6 & 9  & 12 & 15 & 26 & 29 \\
    Green input offset & 1 & 4 & 7 & 10 & 13 & 16 & 19 & 22 \\
    Green selector     & 1 & 4 & 7 & 10 & 13 & 24 & 27 & 30 \\
    Blue input offset  & 2 & 5 & 8 & 11 & 14 & 17 & 20 & 23 \\
    Blue selector      & 2 & 5 & 8 & 11 & 14 & 25 & 28 & 31 \\
    \bottomrule
  \end{tabular}
  \caption{Selectors for the overlapping loads. Selectors 0--15 name input bytes 0--15 in \texttt{\$w0}; selectors 16--31 name input bytes 8--23 in \texttt{\$w1}. A zero-producing control follows every listed selector to form unsigned halfwords.}
  \label{tab:revision-msa-selectors}
\end{table}


The destination has two roles: controls before the shuffle, selected data after
it. The worked code retains immutable controls in \texttt{\$w16} through
\texttt{\$w18} and copies them for every group:
\begin{minted}{mars-idea}
    move.v $w4, $w16
    move.v $w5, $w17
    move.v $w6, $w18
    vshf.b $w4, $w1, $w0
    vshf.b $w5, $w1, $w0
    vshf.b $w6, $w1, $w0
\end{minted}
Afterward, \texttt{\$w4/\$w5/\$w6} contain red, green and blue halfwords. A second
group that reuses these results as controls selects according to component values.
It can fail even though the first group was correct. The control tables are
ordinary private read-only constants; they are not extra caller pixels.

\Needspace{18\baselineskip}
\section{Apply the scalar field calculation in every lane}

The layout now matches the arithmetic. Each lane contains one component between
zero and 255, and each lane index names the same pixel in all three registers.
The packer performs the same shifts as the scalar formula:
\begin{minted}{mars-idea}
    srli.h $w4, $w4, 3
    slli.h $w4, $w4, 11
    srli.h $w5, $w5, 2
    slli.h $w5, $w5, 5
    srli.h $w6, $w6, 3
    or.v $w4, $w4, $w5
    or.v $w4, $w4, $w6
pack_fields:
    st.b $w4, 0($a1)
\end{minted}
\texttt{srli.h} shifts each halfword logically to the right, inserting zero bits;
\texttt{slli.h} shifts each halfword left. Neither moves a component into the
neighboring halfword. The immediate is a bit count. Whole-vector \texttt{or.v}
combines the bit patterns; it does not sum components across lanes.

Lane zero follows the earlier trace: 123 becomes 15 then 30720; 45 becomes 11
then 352; 67 becomes 8. The OR yields \texttt{7968} in that halfword.
The store writes \texttt{68 79} as its first two bytes. Lane one supplies the
next two bytes, and so on. Output order comes from the earlier lane arrangement;
parallel arithmetic does not repair a mistaken shuffle.

The complete source also expands packed pixels for comparison. It repeats bits
according to the stated formulas, using a vector group and a scalar remainder.
Its additional interleaving instructions are not prerequisites for this packing
explanation. Inspect that implementation after deriving its output layout, rather
than treating an unfamiliar shuffle sequence as a formula to memorize.

\section{Carry the group invariant into a scalar tail}

After $g$ complete groups, exactly $8g$ pixels have been converted. The source
pointer is $A+24g$, the destination pointer $D+16g$, and the remaining count is
$n-8g$. All earlier output bytes satisfy the same RGB565 rule. These relations
explain the three loop updates and the next eligibility test. Saving the original
count in \texttt{\$v0} keeps the returned count independent of the shrinking
remaining-count register.

When fewer than eight remain, scalar byte loads read one RGB triple, scalar
shifts form the fields, and two byte stores write the low and high packed bytes.
The scalar tail is part of the algorithm, not an approximation or an error case.
Count nine uses one vector group and one scalar pixel; count seven uses only the
tail. Count zero touches no caller data. For every count, the combined loop
initializes exactly $2n$ output bytes and reads only the selected input bytes.

An implementation may use different group sizes or an equivalent scalar schedule.
Behavioral checks cannot prove that it used vector instructions merely because
its output agrees. Conversely, a particular source spelling is not required for
correctness: some wider shifts can be equivalent when their inputs have enough
zero bits and no surviving bit crosses a field boundary. Establish the value
range and follow those bits before declaring two instruction choices different.
The lab uses exact intermediate observations to investigate this worked schedule,
while the public interface checks accept equivalent implementations.

\section{Pack a different representation independently}

Implement \texttt{pack\_rgb555} in \texttt{rgb555.asm}, supplied in the same public
files directory. RGB888 input stays unchanged, but output now has five bits per
component and a zero top bit:
\[
 Q=2^{10}\lfloor R/8\rfloor+2^5\lfloor G/8\rfloor+\lfloor B/8\rfloor.
\]
Store low byte first. The worked pixel $(123,45,67)$ produces hexadecimal
\texttt{3ca8}, stored \texttt{a8 3c}. White becomes \texttt{7fff}, not
\texttt{ffff}. Reusing the RGB565 green shift or red position would implement a
different format even if some dark pixels happen to agree.

Arguments are source, destination and count in \texttt{\$a0} through
\texttt{\$a2}; count is zero through 32. Return the original count in
\texttt{\$v0}. Input has $3n$ bytes, output $2n$, with disjoint arbitrary byte
alignments and possibly interior full-width pointers. Preserve inputs, guards,
metadata, inactive output, the saved registers, frame pointer and stack pointer.
Only selected output bytes may be written; caller reads may use only selected
input or output. Private constants and private stack storage are permitted.
There is no console I/O. Follow the supplied return continuation.

First establish a scalar baseline and verify the field equation. Then choose a
vector arrangement, prove its bounds and retain a scalar tail. Derive a pixel
that distinguishes your two formats, predict its bytes and inspect the first
incorrect instruction if a check disagrees. The public checks accept equivalent
scalar and vector results; demonstrate your chosen vector strategy separately by
explaining and observing its group and tail. A passing output alone is not that
explanation.
